% =====================================================================
%  R. Nielsen, THE NATURAL LAWS OF SOCIETY
%  In: For Idea and Reality.  A Journal, edited by R. Nielsen, Bj. Bjørnson and Rud. Schmidt.  Year 1871. First Volume (Bind 4).
%  Copenhagen: Otto B. Wroblewsky, Publisher; Printed by G. S. Wibe.  Bind 4, pp. 289-315.  (Samfundets Naturlove, in For Ide og Virkelighed.)
%  ENGLISH TRANSLATION, from transcription.tex (Danish), the text of record.
%  GENERATED by ebook-method/for-idee-1870-73/tr_build.py --write from this template and the translators' parts in
%  ebook-method/for-idee-1870-73/tr/ (TRANSLATOR-BRIEF.md; terms: GLOSSARY.md, its key terms approved by Hans
%  2026-10-13/14).  Conventions: TRANSLATION-PLAYBOOK.md.  Page numbers of the print in the margin.
% =====================================================================
\documentclass[12pt,a4paper,openany]{book}
\usepackage[utf8]{inputenc}
\usepackage[T1]{fontenc}
\usepackage{libertinus}
\usepackage{libertinust1math}
\usepackage{textalpha}
\usepackage[english]{babel}
\usepackage[protrusion=true,expansion=true]{microtype}
\usepackage{setspace}
\usepackage[margin=1.3in]{geometry}
\usepackage{fancyhdr}
\usepackage{hyperref}
\hypersetup{pdftitle={The Natural Laws of Society}, pdfauthor={Rasmus Nielsen}, hidelinks}
\setlength{\headheight}{15pt}
\pagestyle{fancy}
\fancyhf{}
\fancyhead[LE,RO]{\thepage}
\fancyhead[C]{\textit{R. Nielsen: The Natural Laws of Society (1871)}}
\setstretch{1.3}
\usepackage{marginnote}
\renewcommand{\thefootnote}{\arabic{footnote})}   % the print numbers each page's notes 1), 2)
\newcommand{\opage}[1]{\marginnote{\footnotesize\textit{[#1]}}}
\newcommand{\essayhead}[2]{\begin{center}{\large #1}\ifx&#2&\else\\[0.3em]{\normalsize #2}\fi\\[0.5em]\rule{3em}{0.4pt}\end{center}\medskip\noindent}
\newcommand{\signature}[1]{\par\smallskip\hfill #1\hspace*{2em}\par}

\begin{document}

% the title page of Bind 4 (1871-1-for-idee-og-virkelighed.pdf), in English
\begin{titlepage}
\begin{center}
\vspace*{3em}
{\Huge For Idea and Reality.}\\[1em]
\rule{8em}{0.4pt}\\[1em]
{\large A Journal,}\\[0.8em]
{\small edited}\\[0.4em]
{\footnotesize by}\\[0.6em]
{\large R. Nielsen,}\\[0.4em]
{\large Bj. Bjørnson and Rud. Schmidt.}\\[1em]
\rule{8em}{0.4pt}\\[0.6em]
{\large Year 1871. First Volume.}\\[0.6em]
\rule{8em}{0.4pt}\\[4em]
Copenhagen.\\[0.4em]
Otto B. Wroblewsky, Publisher.\\[1.5em]
{\small Printed by G. S. Wibe.}\\[3em]
{\small Translated from the Danish}
\end{center}
\end{titlepage}

\essayhead{The Natural Laws of Society.}{}
% --- p. 289 ---
\setcounter{footnote}{0}%
\opage{289}History has its laws just as fully as nature; but
it is more difficult to point out the strict laws in history,
since they do not pass directly into necessity,
but are laws of the spirit, laws of freedom conditioned in manifold
ways.

If we hold to what is common to all humanity, then the
historical movements always proceed from existing conditions of society
that are bound to certain natural conditions. That
society, although a product of the spirit ---
the family spirit, the patriarchal spirit, and so on --- is
subject to a natural necessity, follows directly from the fact that
the human spirit in all its particular manifestations
has its ground in human nature. Man's higher nature
is, despite the Fall and crimes, a spiritual nature; the spiritual
nature is the nobility of the race. The forefathers, the heroes of legend,
the great mythical figures who open the gates of history
and set the driving wheels of events in motion, are, in the midst of
the shimmering photosphere, the strong radiance of the supernatural
with which a divining, wondering imagination has surrounded
them, nevertheless in a spiritually real sense models for
their descendants; they are so precisely in that the spiritual
% --- p. 290 ---
\setcounter{footnote}{0}%
\opage{290}in their deed and being is so natural, and the natural
in turn so spiritual.

According to natural right, that is: nature's own factual
right, the principle holds that everything possible is permitted,
and only the impossible forbidden. Understood in this way the principle
cannot be applied to historical social development.
In society it is precisely the possible that is forbidden;
for to forbid the impossible is superfluous. Here, in relation
to the demands of right, there is much that happens that
ought not to happen, and much that ought to happen that does not
happen; in human society there is, what there
is not in blind nature, difference and strife between
right and wrong.

None the less the strife everywhere points back to a
spiritual ground in nature, and, when the fundamental forces
are stirred, transforms itself into a palpable, tangible difference between
the natural and the unnatural. To be sure, a state of nature
is by no means a state of right; but in the decisive sense
the demand of right must nevertheless rest on an original claim,
an inner, irrefutable natural need. The just is in the
last instance the natural, and the unjust the unnatural;
right is natural, wrong unnatural; it is these simple
fundamental determinations on which everything depends when
we speak of \emph{the natural laws of society}.

If matters stood in such a way that the fundamental demands
that arise from the nature of society could
involuntarily be felt and perceived by all as demands that
were natural to individuals and immediately satisfied
each one's need: then the transition from
society in its natural state to a developed state of right
would proceed harmoniously, and the historical movements would be like
% --- p. 291 ---
\setcounter{footnote}{0}%
\opage{291}the tones of a symphony. But now there is, as is well known,
a disquieting opposition between what essentially accords
with the nature of society and the immediate, highly
various demands of individuals. Hence upheavals,
discords, shrieking dissonances that are not resolved.
Instead of a spiritual development grounded in nature there now arises
a cutting opposition between the natural
and the artificial, and the artificial becomes in many
ways unnatural. To clarify the natural laws
of society it is necessary to consider the tension
just indicated.

Society has, by virtue of its nature, manifold demands, both
strong and valid; but because of the still greater manifoldness
of individual claims, these social natural demands
come at every moment into conflict with one another.
It is a fundamental demand that there shall
be order in society; but it is also a fundamental demand
that individuals, precisely by virtue of this order,
shall enjoy a certain limited freedom. For the maintenance of good
order the powers that be may apply many
kinds of compulsion and pressure; individuals put up with it and find
it natural for the sake of order. All the well-disposed feel
that there is a consideration that puts the private interests of individuals
in the shade, and that this is the consideration of public,
civic order.

But when this consideration has for a long time asserted itself
one-sidedly, then the machinery of state becomes more and
more artificial, the application of pressure and compulsion more and
more arbitrary, and the condition in the end wholly unnatural.
Society approaches a turning point, when millions
must suffer so that a few favored ones may enjoy;
% --- p. 292 ---
\setcounter{footnote}{0}%
\opage{292}and when this turning point has arrived, when it has really
become palpable and manifest to the crowd how unnatural
the condition is: then one can predict with the same certainty
that the existing order will be burst asunder, as
one can predict that an explosion will occur when
fire gets into gunpowder. To be sure, the necessity with which
an oppressed, maltreated mass of people surges up and
overthrows a throne is unlike the necessity with which
gunpowder explodes, but the unlikeness is not absolute;
the only sufficient ground for a revolution
is an immediate necessity corresponding to irrefutable conditions of life,
a social natural necessity.

The demand of freedom breaks through with necessity
and thereby announces itself as an original demand of society.
But when this fundamental demand, however natural
it may be, is asserted one-sidedly, the unnatural
again emerges. It is the crowd that demands freedom
and demands it with passion, but crowd and passion
share the innate weakness that they are unable
to limit themselves: a rational self-limitation is
as foreign to the wild mob as to inflamed
passion; the natural determination in which they both,
mob and passion, agree is unbridledness.
But for society unbridledness is the most unnatural
of all; unbridledness is the dissolution of society. It
is rational mechanics. The greater the pressure that awoke
the uprising, the more terrible is the raging of the crowd; but
the wilder the crowd rages, the more quickly the reversal
sets in again. All chains are broken, all bonds loosed;
the abyss yawns, and the abomination of desolation appears on
% --- p. 293 ---
\setcounter{footnote}{0}%
\opage{293}every side. But when all is chaos, confusion and
horror, when the underworld opens its maw, and the brazen walls
of Tartarus are burst, what wonder that anguish
runs through the crowd, that terror paralyzes even the foolhardy,
that from every corner people look out expectantly for
a Jupiter with the thunderbolt. So there must be given
a governing power, a power that brings order --- order
at any price! The ordering power comes from without,
yes, from above; but it does not fall directly down from
heaven; it has in the condition itself its omens and prohibitions.
Despotism on such terms gets the natural law
on its side; the fire of freedom consumed itself: a strong despot begins the exercise of power all over again.  % print: fortærede (as printed)

If there were now nothing else to discover in the development of society
than such divergent extremes, history would, despite the changes in
the outward, be only a repetition of monotonous cycles, in which the expert,
sharp-witted observer would for the most part have to be able to determine
the periods, as one determines the seasons.
But just as the spiritual in man is conditioned by the
natural, so in turn the natural is conditioned by the
spiritual. What the race can find natural at one stage of society,
it can find unnatural at another.

Men's natural feelings and natural demands
rest on an original conception of what man essentially
is, what worth he has, and what goal has been set for him in the
world; from the interaction between the ideal conception
and the natural self-esteem in the form of instinct and impulse
there arises the restless striving, the never-resting
longing toward a distant, perhaps unattainable goal; here is
% --- p. 294 ---
\setcounter{footnote}{0}%
\opage{294}the cause of the movement, the fundamental condition of all historical
progress.

It is not simply either the refinement of morals or
the purification of the mind or the subjugation of the passions that
forms the boundary between the pre-Christian and Christian
periods in the history of the world. To be sure, the so-called
Christian civilization stands in many and essential respects high
above the pre-Christian, but these advantages cannot be taken
singly to its credit. It is notorious that the Christian
virtues have as their companions a train of peculiar vices,
vices that heathen antiquity was not in a position
to breed from such a depth of corruption, because its
consciousness did not reach so far; it is notorious that ecclesiastical
faith has not merely provoked unbelief, but has even, through
human guilt, conditioned a superstition so dark and
abominable that the unnaturalness of heathendom in comparison
with it takes on a tinge of involuntary blindness
and excusable coarseness; it is notorious that crimes of the spirit
in every direction grow with the growing consciousness
of spirit and for that reason --- be it said, if
one will, to the honor of Christian civilization --- have never
taken deeper root or extended farther or
risen higher than precisely in the present age.

The natural conception of what it is to be a human being
was in heathen antiquity quite different from
what it later became through the influence of Christianity; it is
the difference in this natural, mark well, \emph{natural}
conception of man's self-worth that marks the
historical line of demarcation. However prominent the opposition
between Oriental despotism and
Greek freedom, it nevertheless finds on closer
% --- p. 295 ---
\setcounter{footnote}{0}%
\opage{295}consideration its resolution in a common foundation. A trait
from the story of Cambyses can make the matter clear.
The king asks once, in an exuberant mood, his friend and counselor
Prexaspes: What do the Persians say of me? ---
Otherwise, lord, answered Prexaspes, they praise you in everything;
only they think that you are too much given to wine. ---
Now you shall judge for yourself, answered the king, whether the Persians
speak truly, or whether they themselves are out of their minds, when they
say such things. If I hit your son, who stands
there in the vestibule, in the middle of the heart, then it is evident that
the Persians speak untruly, and that I am in my right mind.
--- In the same instant he drew the bow and shot at
the boy, and when he fell, he had him cut open, showed
his father that the arrow stuck right in the heart, and said, quite self-satisfied,
to Prexaspes: So it is evident, then, that it is the Persians
who are out of their minds, and not I; have you ever seen in
the whole world a better marksman? --- Prexaspes answered:
A god himself could not shoot better.

At our present stage of society there is certainly no one,
however coarse and uncultivated he may otherwise be, who
would not find this action in the highest degree revolting
and unnatural. One is revolted not merely at Cambyses's
ruthless cruelty; one is revolted if possible still more
at the slavish resignation with which the father,
to escape the son's fate, bows under arbitrariness
and flatters the despot's omnipotence. But the Persians
of that time, with their social consciousness, did
not find such a royal action so revolting or
unnatural. On the contrary! The cultivated --- Prexaspes belonged
after all to the cultivated --- had to find on the whole that such a thing
was in order. Without an unlimited ruler society could
% --- p. 296 ---
\setcounter{footnote}{0}%
\opage{296}simply not exist, and what was then more natural
than that the Mighty One should act on whim
and at every moment do what he had a mind to.

That this, however many exceptions there may be
in particulars, must have been the prevailing
way of looking at things is also evident from the remarkable fact
that Cambyses's entourage expressed a deep sympathy
for himself at the last. Just as he had
armed himself for battle against the false Smerdis, he received
by an accident a wound that threw him on
his deathbed. He now summoned the foremost Persians,
and in order to convince them that the rebel was a wretched
liar, he revealed to them what had hitherto been kept
secret, that he had really murdered his own brother.
Far from repenting, he only lamented himself and bewailed
his sorrowful fate. The bystanders found it
quite natural: ``when the Persians saw their king weep,
they tore their clothes and wept loudly with him.''

That such a way of thinking was revolting to the free
Greek is certain; but that Greek society
could, though in another way and in another form, permit itself
against outstanding individualities actions that were at bottom
just as revolting as those that distinguish Cambyses,
is also certain. The fact is that the Greeks, despite all
their sense for the free, harmonious development of individuality,
yet, when it came to the point, had as little as
the Persians any decisive respect for personality,
any immediate, pervasive impression of man's
original self-worth. The popular religion, the mythical
foundation upon which the whole culture-edifice shining in beauty,
that adorned Hellas so wondrously,
% --- p. 297 ---
\setcounter{footnote}{0}%
\opage{297}was erected, contains a misapprehension of the eternal
self of the individual, a denial of man's personal being.
The gods in Homer constitute, as is well known, a society
of immortals by themselves, exalted above, and separated by an insurmountable
gulf from ``the mortal generations.'' --- ``Brave son
of Tydeus, why do you ask me of my descent? Like
the leaves of the trees, so are the generations of men:
the wind scatters the leaves on the ground, but when the forests
turn green, fresh ones shoot forth from the branches at the
return of spring. So the generations of men, one is born, another
passes away.'' These words of Glaucus to Diomedes are
more than an involuntary outburst of a momentary mood;
they express a fundamental view pervading the whole of
Greek culture, a presupposition beyond which
the natural self-esteem could never reach. This
misapprehension of the self, this barrier of consciousness, none
of the philosophers was in a position to remove; the only one
who seriously tried it was sentenced to death.

Only with Christianity as a world-historical power
does personal self-esteem gain entrance into society.
It is not the relation to God in its inwardness and purity
that is aimed at here; for the effects of this are not directly
outward-going and visible, but inward-going, invisible.
No, it is the higher conception of man's
vocation as a whole and the feeling of human dignity and human right
grounded upon it that is spoken of.
The Gospel's word of man as created in God's image,
when it sounds in a sensuous ear, has something
highly flattering to self-love; this word has worked
as a leaven in the masses; it has set
% --- p. 298 ---
\setcounter{footnote}{0}%
\opage{298}personal right in an ever clearer light; it has brought the
object of self-contemplation near to the eye and enlarged the visual image.

By mixing the eternal with the temporal, letting
the suprasensible step forth in sensible tangibility,
and taking human passions into service,
as much for enthusiasm and action as for repentance and suffering,
the Church in the Middle Ages was able to make itself
the supreme majesty governing state and people. In comparison
with the lower powers of the world it had originally
right on its side and knew how to make its right valid. For an understanding
of its rise, culmination and decline it is
necessary to make a sharp distinction between the Church
as a religious institution with the means of salvation in its keeping
and the Church as a phenomenon of cultural history, as an encroaching
world power. That Catholicism in the first respect
has unfolded ideal forces and in its way developed strong
religious characters shall never be denied; but nonetheless it
stands firm that its high self-esteem as a world power
was of a selfish nature; the fire in its blood, the marrow in
its bones, the driving force in its head and limbs was, even
when it most clothed itself in the semblance of piety and
holiness, insatiable covetousness, unconquerable lust for power,
incarnate egoism. It is this egoism that
holds the papal church bound under the natural law and marks out
its course. It rose like the sun, and it had to rise.
it had rays that dispelled the darkness; it had a
saving, liberating word; it pointed to an exalted
goal; it spoke in its language the cause of humanity, and
its word penetrated into souls. But by selfishly worldly
means to further a supra-worldly goal, to serve heaven
by ruling on earth and to take honor from
% --- p. 299 ---
\setcounter{footnote}{0}%
\opage{299}men in order to spread God's honor, this fundamental endeavor
contained in its germ a corruption that necessarily
had to lead the development from the first natural over into
the unnatural. Early the hierarchical pretensions of
bishops and archbishops began to be a burden to the newly formed
states; early began the scandals at the
Roman court. And yet it was the clergy who took
up the cause of the oppressed, supported woman and freed
the thrall; it was the clergy who in the name of right and morality
stepped forth against brutal power and
set barriers to barbarous arbitrariness; it was, in
one word, the clergy who in the age of the law of the fist brought order,
and the wildly fermenting society needed order.

It is now easy to see that the same conditions that
from the first brought about the rise of the hierarchy had to bring about
its decline in their consequences. If the domineering, selfish
nature had lain merely in the individuals, not in the Church
itself, then the corruption breaking in, however
widely ramified it might for a time have been, could still
always be regarded as something accidental and transitory;
no one could then determine with certainty whether the papal church
might not allow itself to be purified and reformed to the extent that it
once more took hold of the reins of government, once more attained
its old medieval authority over people and
princes. But the matter stands otherwise. The egoism of the hierarchs
and of the hierarchy is the Catholic church-nature's
own inherent egoism. By this, precisely by this
egoism, it won its world dominion, and by this,
precisely by this egoism, it lost it again; the shadow
is inseparable from the light. It is remarkable to see
how the divine governance has kept this
% --- p. 300 ---
\setcounter{footnote}{0}%
\opage{300}power going. The Church seized the Most Holy, set it up
as a goal, but in reality used it as a means.
To use the Most Holy as a means to the attainment of
worldly ends is crime, a crime of the spirit on
the greatest scale. Those who believe in Providence's immediate,
direct intervention in the course of
world events must, when they fix their gaze on the most audacious
among the popes of the Middle Ages, involuntarily pause and ask:
``How is it possible that God could keep silent and thus let
himself be mocked; how is it possible that the Transfigured One could
sit so calmly at the Father's right hand and look on while his
image and his work were distorted, while the lust for power
robbed the keys of heaven and sold forgiveness of sins for money?''
--- It is justice, but a historical justice
that does not strike the eye. The Highest has been put in;
men have freedom to use it; they also have
freedom to misuse it and to take what follows.
Should he who sits at the Father's right hand and follows with
his eye the course of the world have impatiently shattered the papacy
with a stroke of lightning, then he would already in his own time
have had to begin by shattering Caiaphas and annihilating Pilate.

The historical justice shows itself in this, that the divine governance,
without giving up its right at any point, leaves
to the human race the free use of all the forces and
means that are at its disposal; it is a sign of
the immutability of the eternal decrees and of the infinite
superiority of divine wisdom that the tasks
are solved and that everything nevertheless proceeds naturally.

There is in history just as little trace of any divine   % print: ligesa (as printed)
preferential love for Popery as for Muhammadanism:
the reason why the former came to stand higher
% --- p. 301 ---
\setcounter{footnote}{0}%
\opage{301}than the latter must be sought in this, that the capital of life with which
it traded was richer, more fruitful, nobler.

For the development of society on the large scale the principle
holds without circumlocution that \emph{right is might}, and might
right. --- ``And yet we are to believe in a Providence?'' Yes,
precisely for that reason we can believe in a fatherly Providence; for a
stronger proof of Providence's just governance cannot
be given. Let one take good note of what it is that is dealt with here.
Here we are not dealing with society's justice
toward individuals in particulars, but only with the
historical right by which a social order at a given
time forms itself and asserts its continuance. The legal proposition
has two members, of which the one clarifies the other. Here it
is not said that might is without further ado right; here it is said that
might is right, and why? Because \emph{right} is \emph{might}. In this
lies that it is right that makes might into might;
and if this determination alone held, the goal of perfection
would be attained; it would then be impossible for the holders of power
to misuse power. But history is
in the temporal, and the temporal is the imperfect; so
the opposite determination also holds, the expression of imperfection:
\emph{might is right}. The conflict between
the two fundamental determinations has been illustrated in the most grandiose
manner in the history of the papacy.

Here the beginning is made immediately from the presupposition
that right is might. All the human, natural right
with which the Church from the beginning sat, so to speak, in possession,
it had in the course of time ample occasion to unfold. It
is as if the spirit of history blew upon it and drove it
forward, so that it, for the sake of heaven, had to do all
that it was able to for the education of the human race,
% --- p. 302 ---
\setcounter{footnote}{0}%
\opage{302}for the ordering of society, for the advancement of morality and culture.
This is its speed, its historical fair wind;
the hierarchy has the current with it: forward, forward! in wind
and storm, between breakers and reefs the Church's
ship goes under full sail and cleaves the waves of time.

But when the goal was then attained, when all the papal
commands of right --- non obstante qvocunque --- had become commands of might,
what then? Well, then the converse proposition began,
the proposition of temporality: \emph{might is right}, to press
unrestrainedly forward into the foreground of events. ``It was a
pity,'' people say, ``that the church government should thus
degenerate, that the hierarchy did not keep to the old, venerable,
awe-inspiring track.'' An empty sigh of the heart; there is nothing
here to sigh over. The hierarchy degenerated because it had to
degenerate, it could not do otherwise, it was legitimately excused.
As a stone thrown into the air rises, but with steadily
decreasing velocity, until it reaches a point where it
so to speak turns around and begins to fall with increasing
velocity, so it has gone --- on the way of the spirit ---
with the rise and decline of the hierarchy. What gravity is
for the stone, egoism is for the hierarchy. First the
higher has preponderance over the lower; right is might: the movement
goes upward; but little by little the lower gains
preponderance over the higher; might is right: the movement goes
downward. And just as the stone hovering in the air has,
when extraneous influences are left out of account, an equally
long path to run through downward as it has run through
upward, so the fall of the hierarchy in history has described
a path which, in the sections and length of its curvatures,
gives an inverted image of its rise. It did not
come about that the hierarchy had a lesser power as
% --- p. 303 ---
\setcounter{footnote}{0}%
\opage{303}long as it on the whole made might dependent on right,
and afterwards got an immense power when it really began
to defy all considerations of justice and make right
dependent on might. On the contrary! the last turn,
proceeding from the culminating point, is accurately
designated by history as the beginning self-destruction of might.

To a superficial view it looks as if
the papacy succumbed merely because of the many
and various kinds of outward blows and resistance that it gradually
met. That is a one-sided conception; the essential
ground lay in the papacy itself. In order that this
power might be broken, it had to pass step by step through
all the intermediate stages of self-demolition. It was not enough that
the fire of enthusiasm that had inflamed the Crusaders
was extinguished; the papacy got a substitute: from the embers of faith
superstition blazed up more fiercely, and superstition was
quite as serviceable as faith. It was not enough that the pope, after
having felled the medieval emperor and thereby given
the Middle Ages themselves the first dagger-stab, came into political
dependence on France, for the dependence was
paid for only with honor: the real profit the pope kept
for himself: the Babylonian captivity was truly lucrative.
At Avignon the papal court could put up all church offices
and ecclesiastical dignities publicly for auction;
naturally high bids were made at the auctions.
One would have thought that the last hour of the papacy had struck,
when, since there were now two papal chairs, one at Avignon
and one at Rome, one got at once two, then three and more popes,
who mutually condemned and excommunicated one another, for when
the pope condemns the pope, the papacy has surely condemned
itself, and the self-contradiction is evident. But no!
% --- p. 304 ---
\setcounter{footnote}{0}%
\opage{304}This blow, too, could be cushioned. To put an end to the scandal, the free councils --- at Pisa,
Constance and Basel --- tried, by a supreme effort, to
place the Pope under a constitutional guardianship
and to bring about a real reformation in the head and members of the Church. But this step, which furnished a striking
proof of the impossibility that the hierarchically organized
system should ever be able to reform itself
in an evangelical spirit, turned out precisely to the Pope's advantage.
It was therefore not this or that single thing, but the
many-sided, consistent unfolding of the distortion of the rule of the mighty
that gave the ecclesiastical Babylonian tower-building
its foundational shock. Just because power, blind power,
had become right, the pressure was felt in all relations everywhere
and from all sides; states, classes of society, individuals, lay
and learned, found the condition revolting, intolerable, unnatural.
The heavens were overcast, the air close and the horizon
dark; Luther's word was the lightning in the cloud.

As the papal-theocratic universal monarchy
fell apart and sank, the powers of the states
in the various countries of Europe began to gather their strength and to prepare
themselves for a great future. Under the guardianship of the Church
the state could not come to full clarity about its real
right and significance. As papal vassals the princes were
obliged, in all great, decisive matters,
to obey the Pope's orders; a practical, natural conception
and carrying out of real, present state interests
was under these conditions impossible. The tasks were
worldly, human, altogether natural, but the Church
always knew how to substitute something over-worldly, over-human,
supernatural. Under the pretext of ennobling
% --- p. 305 ---
\setcounter{footnote}{0}%
\opage{305}the earthly by an admixture of the heavenly, it
made the heavenly itself into something earthly, and its interference in
everything became intolerable.

When the thunderstorm of the Reformation had cleansed the air,
the state's breathing became freer, not only in the Protestant
but also in the Catholic countries. The state became, as
the phrase goes, conscious of its autonomy (self-legislation).
At first glance this progress seems to be a great
retrogression; for in classical antiquity --- among the Greeks
and the Romans --- the state was already recognized as the
highest real expression of the divine authority
on earth: the state protected by the gods was
in reality the deity itself present. It must not
be forgotten, however, that the state which gathered itself together out of the ferments of the Middle Ages, the state that took over
the power by inheritance from the Church, came to stand on another and higher foundation
than the ancient one; the modern state may, in a natural
and yet ideal sense, be called the reborn state.

To be sure, politics itself, in spite of the devout
mask and the purple rags of edifying phrases
with which in the sixteenth and seventeenth centuries it
so gladly adorned its Machiavellianism, did not
become more Christian than it had already been in the age of Alexander the
Great and Emperor Augustus; but in its mechanical systems of equilibrium and motion
it acquired new masses,
new forces and new conditions to operate with.
Despite all its sins against the humane, the old
Church had in the course of the ages little by little built up so large a
capital, invested, if we may say so, in the supreme guardianship of humanity,
that even the most autocratic
% --- p. 306 ---
\setcounter{footnote}{0}%
\opage{306}authority could no longer withhold with impunity their
right from the wards.

The state neither could nor would thrust the Church from itself; nor
would it, in the manner of antiquity, put itself in the Church's place.
In officially and solemnly recognizing the Church as
a religious, relatively independent institution, subject to the authority of the state only in
everything worldly, it recognized
at the same time the fundamental conception of human life,
with the high ideal representations to which the Church in dogmas,
ceremonies and images gave vital force and sensuous
presence. So the state was compelled to recognize
that in every human individual there dwells an eternal, personal
self; it was compelled to recognize that, besides the
earthly and temporal aims, an
eternal aim is set for the whole race; that after death there is a judgment; that all,
high and low, mighty and humble, the crowned prince
as much as the most wretched beggar, shall one day appear before the one
just Judge, and each receive according to his deeds.
Here is a firm, unshakable foundation. On this foundation
the rule of the mighty, with its political intrigues,
oppressions and acts of violence, might play its game as freely
as it would: a dam had been set against arbitrariness.
Even if the movements went like a rushing stream, a roaring,
swelling river whose waves here and there flooded
the lower banks, drowned the seed and swept away many a
flower and noble plant: there was nevertheless
on both sides a barrier of rock; here was a spiritual
natural barrier. To be sure, the masses in the darkest
spots of society could still with a certain right be said to be in
a condition of rudeness, ignorance and superstition
that yielded nothing to the old heathenism; but,
% --- p. 307 ---
\setcounter{footnote}{0}%
\opage{307}however strange it may sound, the rudeness
itself was nevertheless of a higher and nobler kind. It was a Christian
rudeness, a Christian ignorance, a Christian superstition. But
--- one says --- rudeness, ignorance and superstition are surely
precisely the un-Christian! Quite so. And yet we call
this new barbarism Christian, because it could not possibly
have come forth except on Christian conditions: it could
only come forth in Christendom. It was this un-Christian element
in Christendom that fermented in the masses and set
the peoples in unrest. This un-Christian element in Christendom became a
new driving force, which gave the movements on the large scale a peculiar
direction and impressed its distinctive stamp on the feeling of the tolerable and the intolerable,
of the \emph{natural and the unnatural}.

The gigantic force with which the Revolution at the end of the eighteenth century
broke forth and swept everything
along with it will, from this point of view, prove to be grounded
in a historical natural necessity.

The absolute monarchy, in which the state as state sought
to realize its ideal, was indeed able to hold the
mutually conflicting elements in a certain equilibrium, to subdue dissension
and to simplify the machinery of government by
gathering all the threads of administration in one hand; but it was
not able to recognize the relative independence of the particular institutions,
not able to reconcile the power of the state
with civil liberties and universal human
rights. In the person of the absolute king the state
itself became in a way personal; but instead of reconciling
the legal with the personal, it let the decisions
in the highest instance rest on grace and disgrace; even when
it meant it for the very best, it gave arbitrariness
free play in sympathy and antipathy and set respect of persons
% --- p. 308 ---
\setcounter{footnote}{0}%
\opage{308}above the law. Standing on worldly ground, it
passed itself off, after the model of the hierarchy, as an order instituted
immediately by God himself: it was a royal theocracy
that took the place of the papal. Issuing from the Middle Ages,
in which every single estate had borne its own banner
as high as it could, the estates of society now were,
nobility, clergy, burghers and rural common people --- it was scarcely
possible to call the peasantry an estate --- all without
exception equal before the throne; but they were for all that far
from being equal before the law. It was not enough that the individuals
personally submitted; each estate, even
the nobility, the first, was in relation to the monarch's most exalted
person a subject, deeply bowing, most humbly submissive
personality. There is a self-drive in this personal element that
first glorifies and then kills the monarchy.

The absolute monarchy was as little a repetition
of Roman Caesarism as of Oriental despotism ---
scenes like that with Prexaspes and Cambyses are once
and for all relegated to the unnatural --- the submission, organized
according to principles of right, of the distinct, stepwise
ascending personalities was surely meant precisely to satisfy
society's rational demands upon a well-ordered
state. But in such a system of personal regard for rights
the distinct rights, favors, privileges
and prerogatives were bound to undermine right
as right, universal right. Precedence for the few was wrong
against the many; and on what, in the end, did all precedence have its
support, its proper center of gravity? In the monarch's personal
grace: right was grace, and grace was right; therein
right was dissolved. It is, seen a little more deeply, not the arbitrariness
of those who govern that misuses the demonstrations of grace
% --- p. 309 ---
\setcounter{footnote}{0}%
\opage{309}and thereby makes grace arbitrary. The matter stands
the other way round; grace as a source of right is arbitrariness itself:
it is the absolutely ruling grace that from first to
last, from the highest to the lowest, made the government arbitrary.
It was no accident that the absolute kings
misused grace; no, it was a spiritual natural necessity;
if grace is in itself a misuse of right,
how could it be used without being misused?
It was, strictly speaking, not the monarchs' fault that
the absolute monarchy fell; it fell --- by royal
grace, according to the law of nature: the germ of death lay from conception
in its organism.

It was the inner, revolting discord between right
as grace and right as right that, through countless
excesses, crimes and blunders, brought on the catastrophe
in France. The Revolution of 1789 became an overthrow of the state
in the most radical sense; it led literally
to a bursting of the state as state. When the National Convention
abolished the royal dignity and the Reign of Terror
began, the mass of the people had simply
abolished the state and put its creatures in the state's place; for
a Robespierre is no statesman, a Robespierre cannot
govern; he can only terrorize. In contemplating
such a natural phenomenon it is not sufficient
to dwell on the outer sides, on the multiplicity of the horrors,
on the sea of flames of the passions or merely on the different
eruption craters of the lava streams; one must go down into the
great abyss, down into the inner life, inquire and ask: whence came
all this? What was that depth of glowing
hatred out of which the destruction arose? It was --- a Christian
revolution, a revolution that was possible only in the midst
% --- p. 310 ---
\setcounter{footnote}{0}%
\opage{310}of Christendom. --- That the masses, when the outer pressure becomes
too great and the authorities of the state themselves, in the unrest, act
gropingly and, divided, betray one another, can in
any age whatever be incited to level all
distinction of estates and to intoxicate themselves with fancies of freedom,
equality and brotherhood, shall never be denied. Such a thing is possible
just as well under heathen as under Christian
presuppositions. But the strength of self-feeling, the certainty
of instinct, the clarity of thought and the resoluteness of will
with which the French Revolution laid down the great fundamental law of the ``rights of man''\footnote[1]{``It is a characteristic trait that when the constituent assembly in Paris (1789) discussed the rights, and Mirabeau, impatient that the discussions took too philosophical a turn and dragged on, exclaimed: N’employez pas le mot de droits, mais dites: Dans l’intérêt de tous, il a été déclaré\dots{} the assembly held fast, with the certainty of instinct, to the point that it was precisely rights, that is, express demands of principle, that were to be carried through.'' On Obstacles to and Conditions for Spiritual Life in the Present Age. P.~34.} cannot
be explained without presupposing Christianity.

The men of the Revolution raged against the Church and positive
religion as against all historical traditions in general;
they wanted to burn off the past, clear
the burnt sites and begin afresh with new men,
new gods and a new reckoning of time. But this revolutionary
raging against Christianity is just an eloquent proof
of how deeply that mighty spirit of the age, though in a
selfish, negative, demonic manner, had absorbed and appropriated
the first humane demands of Christianity.

The Revolution's ``freedom and equality'' is a distortion
of the Gospel's freedom and equality; but by this
% --- p. 311 ---
\setcounter{footnote}{0}%
\opage{311}distortion no historical step has been taken at all
toward the abolition of Christianity; rather an unavoidable
demand has been set for a deeper understanding and more inward
appropriation of the essence of Christianity, a demand
that applies just as fully to society as to the individual.

The contradiction between ``the Christianity of the New Testament''
and ``the Christianity of Christendom,'' a contradiction that
S. Kierkegaard has presented in such eloquent, striking
features that his presentation will undoubtedly one day go
round the world, this cutting, spirit-consuming contradiction
demands a twofold resolution, an ethical-historical one in
the human society and one, conditioned by the life of the congregation,
a religious one in the consciences. The latter shall
not be treated in this connection; here we speak only
of the first.

What S. Kierkegaard has called ``the Christianity of Christendom''
is not the Gospel unto salvation, in that he is entirely
right; but when he simply regards it
as a loathsome monster, a nonentity, that ought to be abolished
the sooner the better, he is obviously wrong. The Christianity of Christendom
is, in spite of all distortions, an irresistible
historical factor, a world power, an element in the divine
governance; it does not beg for the approval of the powers that be,
it pays no heed to the arbitrariness of individuals; it
goes, as the stars go, the way of necessity, its
effects are spiritual natural effects. It accompanies
the generations; it submits itself to the spirits of the age; it
leads them, even against their will, sets them their tasks
and shows them the ideals of the future as in a magic
mirror. What sides the tasks set present each time
to reality, and how the visions in the mirror
% --- p. 312 ---
\setcounter{footnote}{0}%
\opage{312}waver depends, according to the occasion of the times, on the faculty that
acts and on the eye that sees. The generation, the society
that with energy and a sense of freedom walks the way of self-limitation
and moderation receives from the Christianity of Christendom
spiritual help, manifold instruction
and a rich blessing. Where, on the other hand, the thought of Christendom
does not purify and ennoble the disposition, it incites it and strengthens
the selfish fancies; where it cannot
act as a blessing, it acts as a curse;
where it cannot be a remedy, it becomes
a poison.

The law according to which the formation of the poison proceeds can
be clearly pointed out. From a revolutionary negative standpoint
criticism hurls at Christianity the accusation
that, by preaching humility and blind
submission to --- grace, it enervates the race and
steals its best powers. Such assaults of thought
always borrow their direction and speed from that which
is attacked. The restless mind unceasingly plagues the understanding
to put, after all, something other and
better in the place of that which it tears down. So
self-feeling shoots upward, the individuals of the masses rise in price, and
the multitude's equal entitlement to all the goods of life sounds
ever louder. It is the despised, repressed and
denied thought of Christendom that reacts; like a subterranean
fire it goes on heaving the social soil
and screwing the demands upon life up into mountains of presumption.
What went on in ancient heathen Rome on the Mons
sacer was only child's play in comparison with what
at the present moment is seen to go on at Montmartre.

It was ingenious of Herodotus to make Mnemosyne's
% --- p. 313 ---
\setcounter{footnote}{0}%
\opage{313}daughters the Muses of his History, for Mnemosyne means
remembrance, and history is essentially remembrance. But the
historical remembrance is of different kinds. To bring out
the natural side we shall here look away from the recollection
of the deeds of the past which historiography
secures, and consider merely the physiological-psychical remembrance,
the unconscious, half-conscious propagation by which the fathers'
dispositions to vices and virtues are inherited
by the children.

History is a chain of connected links;
even if an immediately following one stands in ever so
sharp opposition to its predecessor, there is nevertheless in
the very break a mark that shows how strong the connective tissue
is. The New is born, and the Old dies, it goes
to the ground. But that the Old has gone to the ground means
that it has become the ground from which the New has
arisen. What follows has thus what is past
in its ground; but what it thus has in its ground
it has in an inward manner in its essence, in its inner life,
in natural remembrance. In individual human life
youth is a remembrance of childhood, manhood
of youth, old age of manhood. In history
the same shows itself. Thus Protestantism
is a remembrance of Catholicism, the Revolution a remembrance
of monarchical absolutism; thus the nineteenth
century in way of thinking, morals and conditions is --- not
a repetition, no, on the contrary --- a remembrance of the
eighteenth. Only when one has convinced oneself that in
history there is given an unconscious, naturally necessary remembrance,
and that this remembrance is the inner, invisible bond
between generations and ages, will one recognize
% --- p. 314 ---
\setcounter{footnote}{0}%
\opage{314}that the course of mankind's life, in spite of all apparent
halts, repetitions and retrogressions, is an unceasing
stream, a forward-advancing movement toward a final
goal. It will then be possible to understand that not only the cultivated,
who have knowledge and overview, calculate on a
future for themselves, but that the masses too, in all their rudeness and
ignorance, have historical instincts.

To be sure, an historical earthquake must make the eyes
swim, when the foundation of society suddenly cracks, and
old, long-buried powers are seen to raise their heads
and begin anew; one despairs of right and reason;
one grows dizzy in the circuit and says with the Preacher: ``What
has been, that shall be; what has happened, that shall
happen, and there is nothing new under the sun.'' This is once more
the worn-out history, with the rising of kingdoms and the
fall of kingdoms, with the victory of the one and the defeat of the other, with
the patrician barbarism of orders and the plebeian barbarism of freedom;
where is the progress? When the medieval-Roman
god of the Church and the ancient-Roman god of the state make common cause,
when canon law and the law of the cannon agree that
might is right: what wonder that it falls upon the
powerless!

In order to be able to keep upright in the earthquake and, under
the great general eclipse, to distinguish ghosts from
beings of flesh and blood, it is necessary to
lean on the laws of development and to try to make clear to oneself the
historical remembrance. Papal absolutism is condemned
to death, condemned by the spirit of history. Its time is past;
it is a funerary monument, an image, a shadow
of remembrance: if it comes again, it is only as a revenant.
Royal-imperial absolutism is condemned to
% --- p. 315 ---
\setcounter{footnote}{0}%
\opage{315}death, condemned by history itself; it can rise up,
but not stride forward: it is weighed down by the fetters
of the past. If the phantom returns from its realm of shadows,
it is as a skeleton with a battle-sword for a scythe
and a death's-head under the helmet: the Knight of the Black
Eagle comes not to sow, but to reap. Democratic-revolutionary
absolutism is condemned to death,
and already belongs to the realm of the dead. ``The Red Specter''
on the barricades has no blood in its veins, no breath in
its lungs, it has not the life of presence: seen in the dark
it is a dreadful, blood-dripping phantom, but it is
only a phantom.

To determine the orbits of the heavenly bodies one must
calculate the perturbations. The real orbits are
surely very divergent from those that correspond to the
abstract-universal laws; but these divergences are so far
from casting doubt on the validity of the fundamental laws that they in
reality confirm their fulfillment. The perturbations
to which the phenomena of the great laws of society and
the historical orbits are subject cannot, to be sure,
be calculated mathematically in this way; but it is nonetheless
true that they all work out in the fulfillment of the laws. If there are
many unknown quantities here, there are also many
equations here. The final equation is originally determined by
human nature, constant in its ground, and by the thought of Christ,
unchangeable under all changes.
\signature{R. \emph{Nielsen}.}

\end{document}
